% Target language: Brazilian Portuguese (pt-BR).
% Source: content/en-US.tex, checked against content/zh-CN.tex at upstream commit 0d064f6.
% Translation: AI-assisted; needs native-speaker review before publication.
\StoryParagraph{3}{O peso de um único óvulo: 3,5 microgramas.}

\StoryParagraph{3}{O peso de cinquenta milhões de dólares em espécie: 2,5 toneladas.}

\StoryParagraph{3}{Em uma ligação feita do Montage Laguna Beach, Jing Tian me pediu o segundo, oferecendo o primeiro como garantia.}

\StoryParagraph{3}{Ela tinha ido até lá para fazer uma coleta de óvulos. O bebê, que seria gestado por outra mulher, era esperado para 2027. Se tudo avançasse depressa o bastante, ainda poderia nascer no Ano do Cavalo. A ideia tinha partido dela. Na clínica, disse: sem cinquenta milhões, sem coleta.}

\UnprintedPause


\StoryParagraph{0}{O Montage Laguna Beach fica sobre uma falésia à beira do Pacífico. Reservamos o andar inteiro por trinta dias, ao custo de um milhão de dólares, para preservar sua privacidade. Ela foi embora depois de oito dias.}

\StoryParagraph{0}{Nos vinte e dois dias restantes, as camareiras do hotel continuaram subindo todos os dias, como de costume. Minha assistente me contou. Ela viu isso uma vez, quando voltou para entregar a chave: uma mexicana empurrava um carrinho, entrava em um quarto após o outro, trocava as toalhas, limpava as bancadas, ajeitava os lençóis, saía e fechava cada porta atrás de si. Duas horas para fazer o andar inteiro.}

\StoryParagraph{0}{Minha assistente perguntou por que ela continuava fazendo aquilo se ninguém estava hospedado ali.}

\StoryParagraph{0}{Ela respondeu que, segundo o registro, havia alguém.}

\StoryParagraph{0}{Versalhes, sem sua rainha.}

\StoryParagraph{0}{Foi então que percebi que Jing Tian gostava daquilo. Jing Tian gostava que um lugar permanecesse vazio por causa dela. Que nada acontecesse.}

\StoryParagraph{3}{Além do Montage Laguna Beach, também permanecia vazio um Gulfstream G700, parado sem jeito em Van Nuys enquanto sua tripulação se apresentava para o trabalho todos os dias, como de costume.}

\StoryParagraph{3}{No oitavo dia, ela fez aquela ligação. Depois disso, ninguém voltou a vê-la.}

\NarrativeBreak


\StoryParagraph{0}{Em 2007, alguém já queimava dinheiro por ela. Naquele ano, eu tinha acabado de entrar na Universidade de Pequim.}

\StoryParagraph{3}{É claro que não era eu quem queimava o dinheiro.}

\StoryParagraph{3}{O QQ abriu uma janela com ela como garota-propaganda. Foi naquele velho Lenovo IBM, com uma resolução tão ruim que tudo parecia borrado. Salvei a imagem.}

\StoryParagraph{3}{No dormitório da Universidade de Pequim, onde o aluguel de um ano custava mil yuans, entrei na internet, encontrei-a no Xiaonei e enviei uma solicitação de amizade. O Xiaonei exigia que se acrescentasse uma mensagem. Escrevi e apaguei, apaguei e reescrevi, durante uns quarenta minutos. No fim, enviei: “Olá.”}

\StoryParagraph{3}{Ela não aceitou.}

\StoryParagraph{3}{Talvez nunca tivesse usado o Xiaonei.}

\StoryParagraph{3}{Nos anos que se seguiram, fosse o que fosse que eu fizesse, uma pergunta sempre voltava: até onde precisaria chegar para que ela me aceitasse? Pensei nisso ao alcançar meu primeiro bilhão. Aos dez bilhões, pensei de novo.}

\StoryParagraph{3}{Com o tempo, parei de pensar nisso. O Renren já tinha falido havia anos.}

\StoryParagraph{3}{Nunca enviei uma segunda solicitação.}

\StoryParagraph{3}{Em 2011, fazia frio em Pequim. Comprei um casaco de plumas por cento e cinquenta yuans. Lembro o preço porque hesitei durante três dias.}

\StoryParagraph{3}{Naquele mesmo ano, ela estrelou um filme de cento e cinquenta milhões de yuans, \textit{The Warring States}. Sun Honglei, Francis Ng, Kim Hee-sun, Kiichi Nakai e Jiang Wu tinham sido trazidos como mero pano de fundo, assim como o Montage Laguna Beach e o Gulfstream G700 dezenove anos mais tarde. Eles não faziam ideia de por que estavam ali. Tinham vindo de tão longe; certamente alguma coisa teria de acontecer. Nada aconteceu.}

\StoryParagraph{3}{Depois veio \textit{The Great Wall}, com um orçamento de cento e cinquenta milhões de dólares. Zhang Yimou. Andy Lau. Matt Damon. Depois, Jackie Chan. A lista do universo de Jing Tian não parava de crescer.}

\StoryParagraph{3}{Àquela altura eu já estava nos Estados Unidos: estudava, abria uma empresa, ocupava-me cada vez mais. A única coisa que nunca mudou foi que, sempre que trocava de aparelho, eu transferia aquela imagem do QQ, por mais trabalhoso que fosse.}

\StoryParagraph{3}{Dezenove anos depois, meu nome entrou para a lista.}

\NarrativeBreak

\StoryParagraph{3}{Nós nos conhecemos no Rosewood Hong Kong. Eu sempre imaginara que ficaria nervoso no dia em que a encontrasse.}

\StoryParagraph{3}{Não fiquei.}

\StoryParagraph{3}{Ela era mais magra do que nas fotos, estava vestida sem formalidade e recostada em um sofá. Olhei para ela e, de repente, não consegui lembrar o que exatamente eu esperara por tanto tempo.}

\StoryParagraph{3}{Levei-a para ver a banana de Cattelan. Ela ficou olhando por muito tempo e então perguntou: “Quanto custa isto?”}

\StoryParagraph{3}{Seis milhões e duzentos mil dólares.}

\StoryParagraph{3}{Ela disse: “E quando ela apodrecer?”}

\StoryParagraph{3}{Não respondi. Ela disse: “Deixa para lá. Vamos ver um filme.”}

\StoryParagraph{3}{Eu disse: “Está bem.”}

\StoryParagraph{3}{Ela disse que ninguém além de nós dois poderia ficar na sala.}

\StoryParagraph{3}{Eu disse: “Está bem.”}

\StoryParagraph{3}{Saí para fazer uma ligação. Quando voltei, ela ainda estava parada diante da banana.}

\StoryParagraph{3}{O filme era \textit{Zootopia 2}.}

\StoryParagraph{3}{Uma sala de cinema para duas pessoas. Esse tipo de arranjo costuma ser reservado às sessões em que o Departamento Central de Propaganda avalia um filme.}

\StoryParagraph{3}{Mandei minha assistente comprar todos os lugares do cinema.}

\StoryParagraph{3}{Eu já tinha lido artigos que denunciavam bilheterias falsas: todos os ingressos vendidos, apenas duas pessoas na sala. Agora eu entendia. Talvez nem sempre fosse por causa da bilheteria.}

\StoryParagraph{3}{Comprar todos os lugares não foi difícil. O problema eram os ingressos já vendidos. Meus assistentes e seguranças chegaram cedo e esperaram na entrada. Sempre que aparecia alguém com ingresso, entregavam dinheiro em espécie.}

\StoryParagraph{3}{Um casal chegou com o filho. A criança já estava na porta, segurando um balde de pipoca no alto.}

\StoryParagraph{3}{Uma assistente estendeu um maço de dinheiro. A criança não entendia por que não podiam entrar. Os pais pegaram o dinheiro, agradeceram, não perguntaram nada e foram embora.}

\StoryParagraph{3}{Vinte e seis pessoas aceitaram o dinheiro naquele dia.}

\StoryParagraph{3}{Depois do filme, ela trocou a foto do perfil no WeChat por uma imagem da policial Judy Hopps.}

\StoryParagraph{3}{Naquela noite, coloquei a mão nos peitos dela. Ela me afastou.}

\StoryParagraph{3}{Achei que ela não queria que eu a tocasse. Recolhi a mão e estava prestes a dizer alguma coisa.}

\StoryParagraph{3}{Ela não disse nada. Sentou-se, tirou um pequeno estojo da bolsa e abriu o zíper. Dentro havia uma fileira de objetos, como um minúsculo jogo de ferramentas.}

\StoryParagraph{3}{Ela disse: “Suas unhas me machucam.”}

\StoryParagraph{3}{Olhei. Eu sempre as cortara de qualquer jeito com um cortador de unhas. As bordas ficavam afiadas. Normalmente, eu nem percebia.}

\StoryParagraph{3}{Ela disse: “Venha aqui.”}

\StoryParagraph{3}{Ela pôs minha mão no colo e lixou cada unha, uma por uma.}

\StoryParagraph{3}{Primeiro, o lado áspero. Quando terminou com ele, virou a lixa e depois mudou de lado outra vez. De vez em quando, erguia um dedo contra a luz, olhava e continuava lixando.}

\StoryParagraph{3}{Perguntei: “Quanto tempo isso vai levar?”}

\StoryParagraph{3}{Ela disse: “Até elas pararem de arranhar as pessoas.”}

\StoryParagraph{3}{Dez dedos. Ela levou quase uma hora.}

\StoryParagraph{3}{Enquanto trabalhava, ela me ensinou: este lado dá forma, este lado dá polimento e, com o último, não se deve esfregar para a frente e para trás. É preciso seguir em uma única direção.}

\StoryParagraph{3}{Eu disse: “Está bem.”}

\StoryParagraph{3}{Quando terminou, ela virou minha mão, examinou-a e passou os dedos sobre ela.}

\StoryParagraph{3}{Ela disse: “Pronto.”}

\StoryParagraph{3}{“Pode tocar onde quiser.”}

\StoryParagraph{3}{Em certo momento, ela disse uma coisa. Disse: “Se um dia eu for embora, nunca mais ninguém vai polir suas unhas para você.”}

\StoryParagraph{3}{Sorri. Não disse nada.}

\StoryParagraph{0}{Três meses depois de conhecê-la, mandei alguém montar um dossiê. Os projetos dela ao longo dos anos, parceiros comerciais, participações societárias, Jing Guoqing, Tian Xin'ai, todas as contas relacionadas, o imóvel alugado em Pequim, o imóvel comprado em Xangai e algumas outras coisas. Levou uma semana. Ao todo, quarenta e tantas páginas. Havia coisas sobre homens e mulheres também. Achei que eram o que menos importava.}

\StoryParagraph{0}{Mais tarde, ela me contou algumas das coisas que estavam ali. O que me contou não era exatamente igual ao que estava escrito. Não disse nada.}

\StoryParagraph{0}{Algumas coisas ela nunca me contou. Eu nunca perguntei.}

\StoryParagraph{0}{Acreditava em tudo o que ela dizia.}

\StoryParagraph{0}{Guardei aquele dossiê. Ainda o tenho.}

\StoryParagraph{3}{Mas não acreditava em uma só palavra.}

\NarrativeBreak

\StoryParagraph{3}{Ela disse: “Hong Kong está lotada demais. As pessoas vão me reconhecer. Vamos para algum lugar.”}

\StoryParagraph{3}{Eu disse: “Está bem.”}

\StoryParagraph{3}{Ela apontou Tenerife no mapa.}

\StoryParagraph{3}{Eu disse: “Está bem.”}

\StoryParagraph{3}{Saí para fazer uma ligação. Só quando voltei descobri que o lugar ficava no Atlântico, ao largo da costa noroeste da África. Ao meio-dia do dia seguinte, nosso grupo de trinta pessoas partiu. Um A330 esperava no aeroporto.}

\StoryParagraph{3}{No quarto do avião, ela me perguntou: “Meu nome soa bonito?”}

\StoryParagraph{3}{Eu disse: “Sim.”}

\StoryParagraph{3}{Ela disse: “É a junção do sobrenome do meu pai com o da minha mãe. Não é uma coincidência engraçada?”}

\StoryParagraph{3}{Eu disse: “Gravei seu nome desde o primeiro dia.”}

\StoryParagraph{3}{Ela me abraçou e disse: “Não me chame mais de Jing Tian. Me chame de Mamãe.”}

\StoryParagraph{3}{“Está bem, Mamãe.”}

\StoryParagraph{3}{Quando pousamos, o comandante veio até mim e disse: “Chefe, de início a ilha não queria nos deixar pousar. Nunca tinham visto um jato particular desse tamanho.”}

\StoryParagraph{3}{As regras na ilha eram as mesmas: ninguém podia ficar perto de nós. Meus assistentes cuidaram disso. Ninguém perguntou por quê. Desta vez, eu não perguntei quantas pessoas.}

\StoryParagraph{3}{O Ritz-Carlton dava para o Atlântico, com janelas do chão ao teto grandes demais. Ela não conseguia dormir com o menor vestígio de luz. Os assistentes cobriram o quarto inteiro com fita preta. Do lado de fora, o dia estava claro; dentro, era escuro como uma cela. A fita descolava ao sol e precisava ser reaplicada todos os dias.}

\StoryParagraph{3}{Eram onze assistentes ao todo. Chegavam às 6h30 todas as manhãs e só iam embora depois que adormecíamos. Duas moças cuidavam da fita: uma cobria os caixilhos das janelas; a outra, as luzes indicadoras do ar-condicionado e das tomadas. Quando terminavam, conferiam tudo três vezes. Compraram trinta rolos de fita e depois mais vinte. Descobriram que a fita aplicada depois das três da tarde se soltava com mais facilidade, porque o vidro estava mais quente naquele horário, então passaram a colocá-la pela manhã. Ninguém lhes ensinou isso. Descobriram sozinhas. Discutiam o assunto como se fosse um ofício.}

\StoryParagraph{3}{Ela sabia o nome daquelas duas moças. Eu não.}

\StoryParagraph{3}{Paramos à beira da estrada, no Parque Nacional do Teide, para ver o pôr do sol. Os assistentes tiraram rosas do carro e abriram champanhe.}

\StoryParagraph{3}{Mamãe contemplou o mar de nuvens e disse: “Lindo. Mas não tão lindo quanto o pôr do sol em Delingha, Qinghai.”}

\StoryParagraph{3}{Eu sou de Qinghai.}

\StoryParagraph{3}{Não disse isso.}

\StoryParagraph{3}{À noite, observamos as estrelas do alto da montanha. Vestíamos casacos acolchoados e ainda sentíamos frio.}

\StoryParagraph{3}{Mamãe disse: “Faça um pedido.”}

\StoryParagraph{3}{Perguntei: “O que você pediu?”}

\StoryParagraph{3}{Mamãe disse: “Não vou contar.”}

\StoryParagraph{3}{Mamãe disse: “É tão lindo. Mas faz frio demais. Podemos ir para as Maldivas?”}

\StoryParagraph{3}{Eu disse: “Tudo bem.”}

\StoryParagraph{3}{Fazia quinze graus abaixo de zero em Tenerife e vinte e sete acima nas Maldivas. Nove mil e quinhentos quilômetros em linha reta. Um voo de onze horas. Mas o aeroporto das Maldivas era pequeno demais, então, no meio do caminho, tivemos de trocar nosso jato de fuselagem larga por um avião menor.}

\StoryParagraph{3}{Mamãe disse: “Aviões pequenos fazem muito barulho. Os grandes são melhores.”}

\StoryParagraph{3}{No quarto do avião a caminho das Maldivas, Mamãe me contou que seu pedido tinha sido ter um filho comigo. Teria de ser por gestação de substituição.}

\StoryParagraph{3}{Nunca perguntei por quê. Achei que ela tivesse medo da dor.}

\StoryParagraph{3}{No meio do voo, Mamãe de repente se sentou e procurou o telefone. Conectou-se à Starlink e disse: “Ah, não.”}

\StoryParagraph{3}{Perguntei: “O que foi?”}

\StoryParagraph{3}{Ela disse: “Alguém está me ameaçando. Ele está atrás de mim há três dias e eu não respondi.”}

\StoryParagraph{3}{Eu disse: “Quanto querem?”}

\StoryParagraph{3}{Ela disse: “Você não precisa se envolver.”}

\StoryParagraph{3}{Ela digitou por muito tempo, apagando e mudando, mudando e apagando. Quando terminou, deixou o telefone com a tela virada para baixo sobre o colo. Pouco depois, virou-o para olhar de novo.}

\StoryParagraph{3}{Perguntei quem era.}

\StoryParagraph{0}{Ela disse: “Alguém que trabalhava comigo. Na área financeira. Foi preso. Assumiu a culpa.”}

\StoryParagraph{3}{“Agora ele saiu, não tem para onde ir e insiste em me procurar.”}

\StoryParagraph{3}{Eu disse: “Então chame a polícia e mande prendê-lo.”}

\StoryParagraph{3}{Ela olhou para mim. “Você não entende.”}

\StoryParagraph{3}{Eu disse: “Você pagou a ele?”}

\StoryParagraph{3}{Ela disse: “Sim.”}

\StoryParagraph{3}{“Quanto?”}

\StoryParagraph{3}{“Oitenta mil.”}

\StoryParagraph{3}{Eu disse: “Você levou três dias para pensar nisso?”}

\StoryParagraph{3}{Ela não disse nada.}

\StoryParagraph{3}{Ela mal falou durante o resto do dia. Só adormeceu depois que pousamos.}

\StoryParagraph{0}{A mala de Mamãe foi a primeira coisa a chegar ao Soneva Jani. A bagagem de trinta pessoas não cabia em um único avião pequeno, por isso foi dividida entre três. Mamãe tinha apenas uma mala. Ela ganhou um avião só para si. Mamãe disse para não a colocarem junto com a de mais ninguém.}

\StoryParagraph{0}{O Soneva Jani era famoso por seu lema “no news, no shoes”. Mamãe pôde relaxar.}

\StoryParagraph{3}{Os mordomos do hotel, seguranças e assistentes somavam vinte e seis pessoas, e mais dois operadores de drones nos seguiam à distância. Os assistentes começaram outra vez a cobrir tudo com fita: a luz do ar-condicionado, as luzes das tomadas, repetidas vezes.}

\StoryParagraph{0}{À noite, trocávamos histórias sobre nossos ex como quem troca senhas, como uma dinastia conferindo os nomes de suas eras. Mamãe não tolerava nenhuma divergência. Mas havia uma palavra proibida para ela: não se podia mencionar o atleta.}

\StoryParagraph{0}{Não o mencionei, embora soubesse que a tatuagem nas costas de Mamãe tinha sido feita para o atleta.}

\StoryParagraph{3}{Talvez fosse grande demais para ser apagada.}

\StoryParagraph{0}{Ela disse: “Agora somos uma família. Não foi fácil para meus pais me criarem. Você deve a eles um presente nupcial.”}

\StoryParagraph{0}{Eu disse: “Está bem.”}

\StoryParagraph{0}{Ela me deu dois nomes, Jing Guoqing e Tian Xin'ai, e dois números de conta.}

\StoryParagraph{0}{Trinta milhões. Em renminbis.}

\StoryParagraph{0}{Fiz a transferência. Eu tinha visto aqueles dois nomes pela primeira vez nas quarenta e tantas páginas.}

\StoryParagraph{3}{Pouco depois, ela me entregou o telefone. Tian Xin'ai respondera com uma só palavra: “Recebido.”}

\StoryParagraph{0}{Certa vez, Mamãe disse: “Sabe no que penso quando peço essas coisas?”}

\StoryParagraph{0}{Eu disse: “Em quê?”}

\StoryParagraph{0}{Ela disse: “Fico pensando se um dia você vai dizer não.”}

\StoryParagraph{0}{Eu disse: “Não vou.”}

\StoryParagraph{0}{Ela disse: “Assim não tem graça.”}

\StoryParagraph{3}{Pensei: o que poderia existir que eu não fosse capaz de fazer?}

\StoryParagraph{0}{Quando o tédio ficava grande demais, eu mexia na imagem do QQ de 2007, tentando transferi-la para um aparelho novo. Mamãe se inclinou para olhar. “Por que você guarda às escondidas a foto de outra mulher? Quem é esta?”}

\StoryParagraph{0}{Fiquei calado. Não disse nada.}

\StoryParagraph{0}{A imagem estava borrada demais, então ela saiu para fazer outra coisa.}

\StoryParagraph{0}{Naquela noite, apaguei a imagem.}

\StoryParagraph{3}{No dia seguinte, recuperei-a de um telefone antigo.}

\StoryParagraph{0}{Mamãe também continuava rolando a tela do telefone. Perguntei o que ela estava vendo.}

\StoryParagraph{0}{Ela disse: “Nada. Só toda aquela bobagem de briga pelo primeiro nome nos créditos.”}

\StoryParagraph{0}{A briga pelo primeiro nome: um filme só pode ter uma estrela listada em primeiro lugar. Mesmo entre o protagonista e a protagonista, apenas um pode vir primeiro.}

\StoryParagraph{0}{Eu disse: “Que tédio.”}

\StoryParagraph{0}{Ela disse: “É. E o pior é que às vezes surgem casais no set. Você passa quarenta ou cinquenta dias filmando, quase sempre contracenando com o mesmo homem. Transa à noite e briga pelo primeiro nome nos créditos durante o dia.”}

\StoryParagraph{0}{Eu disse: “Isso é desumano.”}

\StoryParagraph{0}{Ela disse: “É.”}

\StoryParagraph{0}{“Mas, quando as filmagens terminam, não importa quantas vezes vocês tenham transado: agem como se nada tivesse acontecido.”}

\StoryParagraph{3}{Então ela continuou lendo por mais meia hora.}

\StoryParagraph{0}{Um novo ano. Em 2 de janeiro de 2026, no Soneva Jani, pedi Mamãe em casamento.}

\StoryParagraph{0}{Eu havia decidido na manhã de 1º de janeiro. O anel estava em Hong Kong. Não havia onde comprar um na ilha, então uma assistente pegou um hidroavião até Malé. Voltou naquela tarde com o anel dentro de um envelope do hotel.}

\StoryParagraph{0}{Mamãe disse: “O diamante é pequeno demais.”}

\StoryParagraph{0}{Mamãe disse: “Tudo bem. Depois você compensa com um maior.”}

\StoryParagraph{0}{“E você pode vir a Pequim?”}

\StoryParagraph{0}{Fiquei em silêncio.}

\StoryParagraph{0}{Ela achou que eu não queria ir.}

\StoryParagraph{0}{Mamãe disse: “Por que você é tão bom comigo?”}

\StoryParagraph{0}{Eu disse: “Porque gosto de você.”}

\StoryParagraph{0}{Ela disse: “Você era assim com as outras pessoas antes?”}

\StoryParagraph{0}{Eu disse: “Não.”}

\StoryParagraph{3}{Ela sorriu e não disse nada.}

\StoryParagraph{0}{Certa manhã, acordei cedo. Ela não estava na cama.}

\StoryParagraph{0}{Saí e a encontrei sentada no terraço, de frente para o mar. O mar inteiro estava negro. Não havia luzes acesas. O dia ainda não tinha raiado.}

\StoryParagraph{0}{Perguntei por que estava acordada tão cedo.}

\StoryParagraph{0}{Ela disse: “Estou contando os dias.”}

\StoryParagraph{0}{Eu disse: “Contando o quê?”}

\StoryParagraph{0}{Ela disse: “Volte a dormir.”}

\StoryParagraph{0}{Sentei-me ao lado dela por algum tempo, sem saber o que dizer. Depois voltei a dormir.}

\StoryParagraph{3}{Quando acordei de novo, ela já estava maquiada.}

\StoryParagraph{0}{Em catorze dias nas Maldivas, ninguém a reconheceu.}

\StoryParagraph{0}{Na décima terceira noite, no restaurante, ela tirou os óculos escuros. Pouco depois, tirou também o chapéu e o deixou sobre a mesa. Passou a refeição inteira assim. Na mesa ao lado havia um casal alemão idoso. O atendente era do Sri Lanka. Ninguém olhou para ela uma segunda vez.}

\StoryParagraph{0}{Depois do jantar, ela tornou a pôr o chapéu.}

\StoryParagraph{0}{Mamãe disse: “Vamos voltar.”}

\StoryParagraph{0}{Eu disse: “Não deveríamos ficar mais tempo?”}

\StoryParagraph{3}{Ela disse: “É silencioso demais.”}

\UnprintedPause

\StoryParagraph{0}{De volta a Hong Kong, Mamãe conheceu meus pais.}

\StoryParagraph{0}{Naquela noite, ela disse que sua família precisava se preparar para o encontro entre nossas famílias e para o casamento. Eu deveria enviar o máximo possível. Ainda para Jing Guoqing e Tian Xin'ai.}

\StoryParagraph{0}{Não perguntei por quê.}

\StoryParagraph{0}{Eu disse: “O limite diário de cada cartão é de um milhão.”}

\StoryParagraph{0}{Mamãe disse: “Devagar demais.”}

\StoryParagraph{0}{“Use cinco cartões.”}

\StoryParagraph{0}{Cinco cartões, duas contas, alternando entre elas. No terceiro cartão, hesitei. Mesmo assim, concluí todas as transferências.}

\StoryParagraph{3}{Mamãe disse: “Ineficiente demais.”}

\StoryParagraph{0}{Eu disse: “Vamos trazer todos os nossos pais para Hong Kong no Ano-Novo e passar a data juntos.”}

\StoryParagraph{0}{Mamãe disse: “Já combinamos de passar este ano em Hainan. No ano que vem.”}

\StoryParagraph{0}{Durante o Festival da Primavera, os fogos sobre Victoria Harbour atraíram um milhão de pessoas. Ela não estava entre elas.}

\StoryParagraph{0}{Às três da manhã, levei meus pais para casa.}

\StoryParagraph{0}{As ruas estavam vazias. Não havia outro carro esperando no sinal vermelho. Mesmo assim, parei.}

\StoryParagraph{0}{Mamãe disse que não deveríamos nos ver por enquanto. Ela estava prestes a ir a Los Angeles para a coleta de óvulos. Antes, precisava terminar algumas coisas.}

\StoryParagraph{0}{Mamãe perguntou de novo: “Você pode vir a Pequim?”}

\StoryParagraph{0}{Fiquei em silêncio.}

\StoryParagraph{3}{Ela nunca disse o que havia em Pequim. Eu nunca perguntei.}

\StoryParagraph{0}{Mamãe disse que tínhamos de comprar um imóvel em Hong Kong e colocá-lo no nome dela.}

\StoryParagraph{0}{Eu disse: “Tudo bem.”}

\StoryParagraph{3}{Também não nos vimos no Dia de São Valentim.}

\StoryParagraph{0}{Certa vez, ao telefone, ela disse: “Tiraram oito tubos de sangue de mim esta manhã.”}

\StoryParagraph{0}{Eu disse: “Tudo isso?”}

\StoryParagraph{0}{Ela disse: “Foi tranquilo.”}

\StoryParagraph{0}{Eu disse: “Deve ter sido difícil.”}

\StoryParagraph{0}{Depois, falamos de outra coisa.}

\StoryParagraph{0}{Só mais tarde descobri que, para aquele tipo de coleta de sangue, era preciso estar em jejum. Tinha de ser no segundo ou terceiro dia do ciclo menstrual; não se podia errar o dia. Ela provavelmente saiu por volta das seis daquela manhã. Inverno em Pequim. Ainda estava escuro. Depois vieram cada vez mais exames. Eu sabia porque uma nova rubrica apareceu no resumo mensal de pagamentos da minha assistente. Surgiu pela primeira vez no ano anterior. Não era muito, alguns milhares de yuans. Depois disso, aparecia todos os meses, às vezes duas ou três cobranças. Eu as aprovava rapidamente. Valores assim não exigiam que eu visse os detalhes. Só mais tarde soube que eram coletas de sangue, dosagem de seis hormônios, ultrassonografias, AMH. Para saber quanta reserva ovariana ainda restava, quantos óvulos poderiam coletar, se valia a pena continuar.}

\StoryParagraph{0}{Tudo isso foi feito em Pequim.}

\StoryParagraph{3}{Não sei em qual hospital. Não sei se ela foi sozinha ou se alguém a acompanhou.}

\StoryParagraph{0}{Em 25 de fevereiro, ela pousou em Hong Kong.}

\StoryParagraph{0}{Disse que não iria a Los Angeles antes de ver o imóvel em Hong Kong.}

\StoryParagraph{0}{Enviei a ela cerca de doze imóveis que havíamos examinado nos seis meses anteriores. Ela achou todos feios. Eu não tinha visitado nenhum. Meus assistentes haviam passado seis meses percorrendo-os um a um. Eles os enviavam para mim; eu os enviava para ela. Um ficava em Mid-Levels. No vídeo, era possível ver Victoria Harbour. Tinha sido gravado ao entardecer. Por um instante, a câmera oscilou e captou na parede a sombra de quem filmava. Olhei para aquela sombra e pensei: quem é essa pessoa? Quanto caminhou hoje? Quantos imóveis viu? Sabe para quem este está sendo mostrado?}

\StoryParagraph{3}{Havia pouco tempo demais. Ainda assim, ela embarcou no G700 rumo a Los Angeles e voou para o Montage Laguna Beach em 28 de fevereiro.}

\NarrativeBreak

\StoryParagraph{0}{No oitavo dia, Mamãe ligou. Era tarde para ela e madrugada para mim.}

\StoryParagraph{0}{Ela disse: “Cinquenta milhões. Em dólares americanos.”}

\StoryParagraph{0}{Eu disse: “Deixe-me pensar.”}

\StoryParagraph{0}{Ela não disse nada. Eu conseguia ouvir o oceano onde ela estava.}

\StoryParagraph{0}{Depois de algum tempo, ela disse: “Está bem.”}

\StoryParagraph{0}{Ela desligou.}

\StoryParagraph{3}{Durante todo o tempo em que a conheci, eu nunca tinha dito “deixe-me pensar”.}

\StoryParagraph{0}{Usei o Claude Code para consultar as APIs e calcular todas as minhas reservas em dinheiro. A conclusão foi que aquilo não teria efeito algum sobre mim. Conferi de novo. Digitei a pergunta no Claude Code. Claude disse: “Não dê a ela os cinquenta milhões de dólares.”}

\StoryParagraph{3}{Perguntei: “Ela não me ama mais?”}

\StoryParagraph{3}{Claude disse: “No que diz respeito ao amor, não me importo nem entendo. Mas você não pode dar a ela os cinquenta milhões de dólares.”}

\StoryParagraph{3}{Eu nunca tinha recusado nada a ela. Não por generosidade. Por medo. Medo da resposta que viria no segundo seguinte.}

\StoryParagraph{0}{E agora alguém tinha dito aquilo por mim.}

\StoryParagraph{0}{Fiquei olhando para aquela linha por muito tempo. O que me aterrorizava não era Claude ser inteligente. Era Claude não hesitar nem por um instante.}

\StoryParagraph{0}{Não havia conflito.}

\StoryParagraph{0}{Pensei que aquilo devia ser alguma forma superior de existência. Os seres humanos não eram capazes do mesmo.}

\StoryParagraph{3}{Então ela não me amava mais.}

\StoryParagraph{3}{No dia seguinte, ela ligou de novo.}

\StoryParagraph{3}{Mamãe disse: “O oceano é tão grande.”}

\StoryParagraph{3}{Eu disse: “É?”}

\StoryParagraph{3}{Ela disse: “Este andar é grande demais. Dá eco. Alguém veio limpar esta manhã. Eu me escondi no banheiro e esperei a pessoa ir embora.”}

\StoryParagraph{3}{Eu disse: “Posso pedir que não venham.”}

\StoryParagraph{3}{Ela disse: “Não precisa.”}

\StoryParagraph{3}{Ela disse: “A clínica perguntou quando o pai do bebê chegaria.”}

\StoryParagraph{3}{Eu disse: “Posso mandar alguém até aí.”}

\StoryParagraph{3}{Ela disse: “Não precisa.”}

\StoryParagraph{3}{Então perguntou: “Você já decidiu?”}

\StoryParagraph{3}{Não falei. Não porque estivesse pensando. Eu não pensava em nada. Ao telefone, havia o oceano do lado dela e Victoria Harbour do meu. Nenhum dos dois falava. Contei. Onze segundos.}

\StoryParagraph{3}{Jing Tian disse: “Tudo bem.”}

\StoryParagraph{3}{Jing Tian disse: “Eu entendo.”}

\StoryParagraph{3}{Ela desligou.}

\StoryParagraph{3}{Inúmeras vezes tentei lembrar o que mais ela disse naqueles dois dias, além de cinquenta milhões. O oceano é tão grande. Este andar é grande demais. Ela se escondeu no banheiro. Foi só isso. Lembro cada palavra. O que não consigo lembrar é a voz dela---o tom que usou ao dizer: “Eu entendo.”}

\StoryParagraph{0}{Não consigo lembrar. Por fim, parei de tentar.}

\StoryParagraph{0}{Coloquei o telefone com a tela virada para baixo na mesa de cabeceira. Então fiz uma coisa que ainda não compreendo: pus a mão no peito e contei.}

\StoryParagraph{0}{Meu coração batia em ritmo constante.}

\StoryParagraph{0}{Contei de novo. Continuava constante. Pouco mais de sessenta batimentos por minuto, como sempre. Fiquei sentado ali, com a mão pressionada contra o peito, como um médico examinando outra pessoa.}

\StoryParagraph{3}{Pensei: então é assim.}

\StoryParagraph{0}{Não fiz nada durante o resto daquele dia. Fiquei no quarto até escurecer e clarear de novo, observando Victoria Harbour passar do azul ao cinza e ao negro. Às seis da manhã, as torres financeiras do outro lado do porto começaram a se iluminar uma a uma. Uma assistente bateu à porta uma vez para perguntar o que deveria providenciar para o jantar. Eu disse que não precisava.}

\StoryParagraph{3}{Não chorei. Continuei esperando que o choro viesse.}

\NarrativeBreak

\StoryParagraph{0}{Mais tarde, quando minha assistente examinou as despesas comigo, havia um item de indenização ao hotel em Tenerife. Depois que fomos embora, as moças continuaram colocando a fita por mais alguns dias. Ninguém lhes disse que podiam parar. O hotel afirmou que ela havia sido aplicada por tempo demais, vezes demais. O resíduo no caixilho da janela não saía. Foi preciso substituir todo o painel de vidro.}

\StoryParagraph{0}{Eu disse: “Aprovado.”}

\StoryParagraph{3}{Aquele painel de vidro dava para o Atlântico.}

\StoryParagraph{0}{O combustível da aeronave na fatura custara mais do que o andar inteiro no Montage. Perguntei a respeito. Meu assistente explicou que o A330 fora projetado para duzentas pessoas. Com apenas trinta a bordo, ficava leve demais, e o equilíbrio não era adequado para a decolagem. Por isso, a cada vez, precisavam acrescentar combustível até que ficasse pesado o bastante. Quando chegávamos, o combustível extra não podia permanecer a bordo. Tinha de ser despejado.}

\StoryParagraph{0}{Eu disse: “Todas as vezes?”}

\StoryParagraph{3}{Ele disse: “Todas as vezes.”}

\StoryParagraph{0}{Havia também uma caixa enviada de volta pela clínica de Los Angeles. Oito dias de injeções para estimulação ovariana. Precisavam ser mantidas sob refrigeração. Nenhuma tinha sido aberta.}

\StoryParagraph{3}{Havia ainda o depósito para a agência de gestação por substituição: duzentos mil dólares, não reembolsáveis. A mulher que gestaria a criança tinha sido escolhida em janeiro. Era da Califórnia, tinha trinta e quatro anos e dois filhos. A fotografia dela estava no dossiê. Nunca olhei. Ela estava pronta desde janeiro.}

\StoryParagraph{3}{Depois que Jing Tian foi embora, resgatei a fotografia de 2007 e olhei para ela uma vez. Dezenove anos haviam passado. Os pixels de cada aparelho ficavam cada vez melhores, enquanto os daquela imagem pareciam piorar. O rosto dela era um borrão. Ampliei a foto e fiquei olhando por muito tempo, tentando encontrar alguma diferença entre quem ela era naquela época e quem se tornara depois. Não encontrei nenhuma. A fotografia nunca fora nítida o bastante para revelar coisa alguma.}

\StoryParagraph{0}{Um dia, cortei as unhas. Quando terminei, passei um dedo sobre elas. Estavam afiadas.}

\StoryParagraph{0}{Lembrei que aquela coisa ainda estava na gaveta, a que ela tinha deixado para trás.}

\StoryParagraph{0}{Tirei-a de lá e a passei duas vezes pelas unhas. Não parecia certo. Ela dissera que era preciso seguir em uma única direção. Eu tinha esquecido qual.}

\StoryParagraph{0}{Guardei-a de novo.}

\StoryParagraph{3}{Minhas unhas ainda arranham as pessoas.}

\StoryParagraph{0}{Certa manhã, apareceu um lembrete no meu telefone. Minha assistente o havia programado no ano anterior, ao organizar a agenda: três meses antes da data prevista para o parto, preparar o quarto.}

\StoryParagraph{3}{Deslizei o lembrete para dispensá-lo.}

\StoryParagraph{0}{Desde que ela foi embora, não paro de pensar em uma pergunta.}

\StoryParagraph{0}{Se eu tivesse dado o dinheiro a ela naquele dia, teria ficado?}

\StoryParagraph{0}{Pensei nisso por muito tempo.}

\StoryParagraph{3}{Mais tarde, peguei o dossiê e o li de novo. Quarenta e tantas páginas.}

\StoryParagraph{0}{Então me lembrei dos fogos de artifício sobre Victoria Harbour na noite do Festival da Primavera, vistos do último andar.}

\StoryParagraph{0}{Um milhão de pessoas avançou como a maré. Os fogos terminaram. Um milhão de pessoas recuou.}

\StoryParagraph{0}{Um milhão de pessoas.}

\StoryParagraph{0}{Todo ano é igual.}

\StoryParagraph{0}{Hong Kong às três da manhã---}

\StoryParagraph{3}{Nada aconteceu.}
